\section{Introduction}\label{ch2:sec:introduction}


Federal Reserve announcements and communications are associated with
economically meaningful movements in interest-rate and other asset markets
(\cite{vergote2012}; \cite{rosa2013financial};
\cite{jubinski2013fomc}). U.S. Treasury futures trade nearly continuously,
which permits high-frequency research designs around scheduled policy
releases. The empirical association diagnostic in this chapter, however, uses
release-day closes and does not isolate an instantaneous announcement-window
response.

A related text-as-data literature documents associations between textual
sentiment and financial outcomes (\cite{bollen2011twitter};
\cite{ding2015deep}; \cite{ding2016knowledge}; \cite{xu2018stock}). Such
associations do not establish that sentiment is a single causal driver of
market movements. Central-bank communication is multidimensional, and a
document can convey policy-path or economic-information news that a scalar
tone score does not recover.

Policy meetings and institutional documents also form a sparse historical
record, motivating controlled synthetic-data exercises. Synthetic financial
data have been used for robustness testing and model development
(\cite{cao2022synthetic}; \cite{carvajal2022synthetic}), while pretrained
language models can be adapted to generate task-specific examples or
structured text (\cite{schick2021generating}; \cite{peng2023instruction};
\cite{yu2024large}; \cite{li2021synthetic};
\cite{meoni2024generating}). These applications motivate an empirical test;
they do not make an LLM an intrinsically valid institutional simulator.

Accordingly, this chapter treats synthetic text as a conditional research
artifact whose source preservation, target alignment, and delivery must be
measured separately. The implemented Minutes task rewrites a supplied analysis
as one formal paragraph. It does not generate arbitrary counterfactual
scenarios, reproduce the FOMC's full deliberative process, or establish an
improvement in a downstream volatility model.

This chapter studies a shared-parent, two-branch post-training design. First,
\textit{Model chk-1} maps point-in-time macro-financial evidence into analytical
narratives and provides the common analytical representation. From this parent,
\textit{Model chk-2} (repository artifact \texttt{chk3}, checkpoint 318) is
trained to rewrite a supplied analysis as source-preserving
FOMC-Minutes-style prose. \textit{Model chk-3} (repository artifact
\texttt{chk4}) is trained separately to map a target-neutral pre-meeting
analysis into a monetary-policy action: decision SFT checkpoint 38 initializes
GRPO, whose completed run terminates at step 468. Checkpoint 450 is retained
only as the best-observed exploratory post-GRPO evaluation candidate. The
formal selection record has \texttt{selected\_checkpoint\_step=null}, and the
test verdict is \texttt{quality\_not\_demonstrated}; checkpoint 450 is not a
selected, final, or promoted model. Checkpoint 38 is therefore the retained
pre-GRPO state of \textit{Model chk-3}, not a separate paper-level model.

The two downstream branches are evaluated independently. In particular,
\textit{Model chk-3} does not consume the synthetic Minutes produced by
\textit{Model chk-2}. Accordingly, the decision-prediction results are not
interpreted as evidence of the downstream decision utility of synthetic
Minutes, and this chapter does not claim to validate an end-to-end
indicators-to-Minutes-to-decision system. The completed empirical assessments
instead distinguish source-preserving paragraph rewriting from policy-action
performance in the separate decision branch. The Core8 and
sentiment--Treasury records are retained as legacy descriptive evidence, while
the direct sentiment-score closeness diagnostic remains planned rather than a
reported result. The principal decision evidence is the 13-meeting
training-held-out test. The all-237 comparison is a within-release diagnostic,
not a generalization estimate, because it contains 211 training meetings, 13
validation meetings, and the same 13 test meetings. The test has now been
opened and cannot serve as a new prospective sealed evaluation.



\subsection*{Research Motivation}

Monetary-policy prediction remains relevant for asset pricing and risk
management, and text can contain contextual information not summarized by a
small quantitative predictor set. The information clock is nevertheless
decisive. Official Minutes and post-decision statements cannot serve as public
pre-meeting inputs for the same decision they document. As discussed in
Section~\ref{ch2:sec:lr:realtime}, the decision branch is therefore framed as
a historical classification exercise using a target-neutral pre-meeting
brief, not as a claim that released Minutes forecast the contemporaneous
action in real time.


LLMs provide a practical route for connecting structured indicators with
formal policy narratives. This chapter focuses on two narrower and separately
identified downstream tasks built on a common analytical parent. The first
tests whether Minutes-specific SFT improves faithful analysis-to-prose
rewriting. The second studies policy-action prediction through a comparison
between a decision-SFT checkpoint and its post-GRPO descendant. The GRPO score
is treated as an engineered, format-sensitive proxy for action-label agreement
and native response delivery, not as a measure of reasoning quality or overall
decision quality. This branch does not treat the generated Minutes as a
decision-model input. The Core8 intervention, direct sentiment-score closeness,
and historical sentiment--Treasury exercises are supporting diagnostics rather
than evidence that the two branches form a serial decision-support pipeline.


\subsection*{Research Questions}

This chapter investigates two parallel downstream uses of a shared
macro-financial analysis. The primary research question for the
paragraph-rewriting branch is:

\textit{Relative to its analytical parent, does Minutes-specific supervised
fine-tuning improve source preservation and target alignment when rewriting a
supplied macro-financial analysis as FOMC-Minutes-style prose?}

A separate extension compares structured policy-action delivery and decision
quality between the decision-SFT checkpoint and its post-GRPO descendant. The
completed evidence and explicitly identified follow-up diagnostics are
organized around the following questions:

\begin{itemize}
\item[1.] Does \textit{Model chk-2} improve exact source preservation and
reliability-adjusted semantic alignment relative to its immediate
\textit{Model chk-1} parent?

\item[2.] What do the surviving legacy deletion aggregates suggest about
target-relative sensitivity when one supplied Core8 topic block is removed,
given that row-level pairing and seed provenance cannot be reconstructed?

\item[3.] In a future frozen run, how closely do \textit{Model chk-2}'s
meeting-level sentiment scores track the deterministic synthetic reference,
and is it closer than \textit{Model chk-1} and \textit{Model chk-0} under
prespecified score-space metrics?

\item[4.] What does the surviving legacy sentiment--Treasury record document,
and is its provenance sufficient to estimate cross-artifact coefficient
differences?

\item[5.] In the independent decision branch, does the best-observed
exploratory post-GRPO checkpoint exhibit higher native structured delivery than
its decision-SFT initialization without lower direction accuracy, balanced
accuracy, macro F1, or exact-action accuracy?
\end{itemize}

RQ1 and RQ5 have current frozen evaluation evidence. RQ2 is answered only at
the level of legacy deletion aggregates; the paired deletion--neutralization
design remains to be rerun. RQ3 is a planned diagnostic with no active result
in this chapter, and the surviving RQ4 record is insufficient for a verified
model-difference test. None of these questions treats an emitted rationale as
evidence of interpretable or economically valid reasoning, because the
present evaluation contains no independent reasoning-quality measure.


\subsection*{Contribution}

This study makes several important contributions to the literature at the intersection of central banking, financial economics, and artificial intelligence:

\begin{itemize}
\item[] \textbf{Shared-Parent, Two-Branch Design:} We develop a common
analytical model followed by two independently evaluated downstream branches:
analysis-to-FOMC-Minutes-style paragraph rewriting and analysis-to-policy-action
prediction. The decision branch is a separate post-training comparison, not a
validation of the downstream usefulness of synthetic Minutes; its
post-GRPO comparison is descriptive rather than a causal estimate of GRPO.

\item[] \textbf{Source-Preserving Text Rewriting:} We evaluate whether
Minutes-specific SFT improves exact numeric preservation and
reliability-adjusted target alignment when rewriting supplied analyses as
formal FOMC-Minutes-style paragraphs. On the external 128-meeting evaluation
with ten stochastic replicates per model and meeting, the Minutes branch shows
supported numeric and semantic increments over its analytical parent. One
additional date failure and one additional degeneration failure nevertheless
trigger the preregistered zero-tolerance verdict
\texttt{random\_decoding\_regression}, precluding an unconditional improvement
claim.

\item[] \textbf{Independent Decision-Branch Diagnostic:} We compare decision
SFT with its post-GRPO descendant under a common action and response contract,
using the 13-meeting training-held-out test as the principal decision
assessment. We report the 237-meeting comparison only as a within-release
diagnostic because it comprises 211 training, 13 validation, and 13 test
meetings. Policy-action quality, native structured delivery, and the engineered
proxy score are reported separately.


\item[] \textbf{Extension of LLM Applications:} This study applies
domain-specific LLM post-training to controlled institutional-text rewriting
and policy-action diagnostics, extending its use beyond conventional text
classification and information retrieval.

\item[] \textbf{Proxy-Objective Diagnostic:} The evaluated post-GRPO
checkpoint has higher native JSON delivery and a higher replayed proxy score on
the all-237 diagnostic, but checkpoint 38 has higher point estimates on all
four primary policy-decision metrics. On the 13-meeting test, checkpoints 38
and 450 have the same direction and exact-action accuracy and neither produces
a delivery-valid output. No GRPO candidate passes all selection gates, the
formal selected checkpoint is null, and the evidence therefore supports a
decision-quality--delivery trade-off rather than an overall GRPO improvement.

\item[] \textbf{Practical Implications:} The findings identify both the
potential and the current reliability limits of domain-adapted LLMs for
source-controlled policy-text generation and structured policy-action tasks.

\item[] \textbf{Evidence-Layered Evaluation:} We combine the formal
Minutes-fidelity evaluation with legacy descriptive Core8 deletion and
historical sentiment--Treasury records, while identifying the canonical
paired intervention and direct sentiment-score closeness test as unfinished.
Policy-action performance is evaluated separately in the parallel decision
branch; these evidence layers do not constitute a serial end-to-end pipeline.
\end{itemize}


\subsection*{Chapter Structure}

This chapter is organized as follows. Section~\ref{ch2:sec:lr} reviews the
literatures on central-bank communication, real-time policy prediction,
controlled data-to-text generation, domain adaptation, faithfulness
evaluation, and reward-based post-training. Section~\ref{ch2:sec:training-plan}
defines the shared-parent, two-branch post-training design and the
paper-to-repository model mapping. Section~\ref{ch2:sec:the-dataset} describes
the dataset construction process and prompt templates.
Section~\ref{ch2:sec:base_model} introduces the base model, while
Sections~\ref{ch2:sec:training} and~\ref{ch2:sec:parameter} describe the
fine-tuning method and hyperparameter settings.
Section~\ref{ch2:sec:application} presents the evaluation dimensions and
empirical tests, and Section~\ref{ch2:sec:result} reports the training and
evaluation results. Finally, Section~\ref{ch2:sec:conclusion} concludes the
chapter and discusses future research.
